\subsection*{S10. Duplicate keys}

An object carrying the same key more than once.

\begin{lstlisting}[style=json]
{ "a": 1, "a": 2 }
\end{lstlisting}

RFC~8259 says member names \emph{should} be unique and declares the behaviour
undefined when they are not. Implementations genuinely disagree: some keep the
last binding, some the first, some raise an error.

\options
\begin{enumerate}[label=\textbf{(\alph*)}]
  \item \textbf{Reject}, with an error naming the path and the repeated key.
        \selected
  \item \textbf{Last binding wins}, as JavaScript's \texttt{JSON.parse} does.
  \item \textbf{First binding wins.}
  \item \textbf{Treat the key as repeated}, lowering the path to a collection
        as though the document had supplied an array.
\end{enumerate}

\decision{Reject. A document containing a repeated key within any object, at
any depth, is an error naming the path and the key.}

\rationale{The goal is agreement between querying the JSON and querying the
database built from it. Here the JSON side has no defined answer: two
conforming JSON query engines may legitimately disagree about the value of
\texttt{.a} in the document above. Agreement cannot be proved against a
reference semantics that is itself undefined, so admitting duplicates would
leave a hole in the statement at precisely the point where the statement is
the product.

Options (b) and (c) discard data silently. Option (d) preserves it but invents
a collection the document did not ask for, and makes the schema depend on
whether a key happened to repeat. Rejection also costs little in practice:
duplicates are rare, and when they occur they usually indicate a fault in
whatever generated the document --- something the user would rather be told
about than have quietly resolved.}

\implnote{The parser will not perform this check for us. Yojson represents an
object as an association list and preserves both members: \texttt{\{"a":1,
"a":2\}} parses to a two-element list and re-serializes with both, verified
against yojson 3.0.0. The check must therefore be explicit, and applied to
every object at every depth. It is a per-document check, not a corpus-level
one.}

\limitation{Documents that other tools accept will be rejected. Since the
check is per-document, one malformed document rejects the whole input unless
the interface offers to skip it.}
